% ECONOMICS_PROBLEM: {"id": "Q51-363", "title": "Nonmarket climate losses", "jel_code": "Q51", "source_id": "OP-0432"}
% !TeX program = xelatex
\documentclass[11pt,a4paper]{article}
\usepackage[margin=23mm]{geometry}
\usepackage{fontspec}
\setmainfont{Latin Modern Roman}
\usepackage{amsmath,amssymb,mathtools}
\usepackage{xcolor,enumitem,booktabs,longtable,array,xurl,hyperref}
\definecolor{muted}{HTML}{53636C}
\hypersetup{colorlinks=true,urlcolor=blue,linkcolor=blue}
\setlength{\parindent}{0pt}
\setlength{\parskip}{5pt}
\newcommand{\R}{\mathbb R}
\newcommand{\E}{\mathbb E}
\newcommand{\N}{\mathbb N}
\newcommand{\ind}{\mathbf 1}
\newcommand{\Var}{\operatorname{Var}}
\newcommand{\Cov}{\operatorname{Cov}}
\newcommand{\supp}{\operatorname{supp}}
\DeclareMathOperator*{\argmin}{arg\,min}
\DeclareMathOperator*{\argmax}{arg\,max}
\newcommand{\entrymeta}[3]{{\sffamily\small\color{muted}\textbf{\texttt{#1}}\quad|\quad #2\par #3\par}}
\newcommand{\Problemref}[1]{\texttt{#1}}
\newcommand{\JELref}[2]{\texttt{#2} (JEL \texttt{#1})}
\begin{document}
\section*{Nonmarket climate losses}
\textbf{Problem ID:} \texttt{Q51-363}\quad \textbf{JEL:} \texttt{Q51}\par
% BEGIN ECONOMICS STATEMENT
\label{problem:OP-0432}
{\sffamily\small\color{muted}Original subject group: Climate, environment and energy\par}
\label{definitions:problem:30007714}
\textbf{Audit classification:} Background; proposed extension. Added compensation existence/range assumptions and distinguished utility discounting from financial present value.
\subsection*{Definitions and precise research target}
\textit{Editorial formalization.} Consider coastal households in one prespecified country facing two feasible climate-policy paths over thirty years. Let \(p\in\{0,1\}\) index the paths, with \(1\) reducing warming relative to \(0\). For household \(i\), define consumption \(c_{it}(p)>0\), health \(h_{it}(p)\), accessible ecosystem services \(e_{it}(p)\), and an indicator \(d_{it}(p)\) for involuntary relocation in year \(t\). Let lifetime utility be
\[U_i(p)=E\!\left[\sum_{t=0}^{30}\beta^t u_i(c_{it}(p),h_{it}(p),e_{it}(p),d_{it}(p))\right],\]
where \(\beta\in(0,1)\), preferences are stable over the comparison, and the expectation covers specified climate and income risks. Fix a transfer schedule \(\ell_t\geq0\) with \(\sum_{t=0}^{30}\beta^t\ell_t=1\). Define compensating variation \(v_i\) by replacing \(c_{it}(0)\) with \(c_{it}(0)+v_i\ell_t\) until utility equals \(U_i(1)\), when a unique feasible solution exists. The schedule fixes transfer timing under a declared utility-discount normalization; financial present value requires separate discount prices. The research task is to identify population-weighted \(E[v_i]\) and its distribution without counting health, displacement, or ecosystem losses twice through market expenditure measures. Specify revealed-preference or stated-preference assumptions, treatment of cultural losses that cannot be traded, and sensitivity to adaptation and income-dependent valuation.

For a fixed baseline household population, let \(U_i^0(v)\) be the baseline expected utility after the stated consumption-path transfer \(v\ell_t\). Feasible transfers form \(\mathcal V_i=\{v:c_{it}(0)+v\ell_t>0\ \text{almost surely for all }t\}\). Assume \(U_i^0\) continuous and strictly increasing on \(\mathcal V_i\), and \(U_i(1)\) belongs to its range. Then \(v_i=(U_i^0)^{-1}(U_i(1))\) is well-defined; otherwise report nonexistence or a one-sided bound rather than a finite monetary value. \(\ell\) fixes transfer timing, while its displayed normalization uses the utility discount factor by convention; actual financial present value instead requires a stated bond-price vector. Use fixed population weights and \(E|v_i|<\infty\). Nonmarket dimensions enter utility directly and are not added again as separate damage dollars. All expectations use a stated baseline population and probability law. Identification means every admissible data-generating model with the same observed-data law yields the same displayed target; otherwise report the identified set. Exact equations, horizons and assumptions are editorial, rather than source-given canonical conjectures.
\subsection*{Variants and assumption changes}
\begin{enumerate}
\item \textbf{Editorial equivalent-variation variant.} Define \(e_i\) by \(U_i^1(-e_i)=U_i(0)\) using the policy-1 transfer schedule. Compare \(e_i\) with baseline compensation \(v_i\); income effects can prevent equality.
\item \textbf{Editorial nonmonetary variant.} Drop compensability or range inclusion for cultural losses. Report joint health, ecosystem and displacement outcomes under explicitly weighted utility rather than forcing a finite dollar compensator.
\end{enumerate}
\subsection*{References and source audit}
\textbf{1.} SCC synthesis discusses damage persistence, tipping, and nonmarket substitution; proposed estimands are editorial. Locator: Introduction; Section 1.1; Section 2. Full text or primary publication page inspected. URL: \url{https://pmc.ncbi.nlm.nih.gov/articles/PMC11670083/}.

\textbf{2.} Existing biodiversity production model; no exact replacement-frontier open conjecture verified. Locator: Working-paper record and abstract. Primary indexed metadata and abstract checked; direct page and PDF retrieval failed. URL: \url{https://www.nber.org/papers/w32678}.
\begin{enumerate}\raggedright
\item\hypertarget{definitions:source:30007714:1}{} Frances C. Moore; Moritz A. Drupp; James Rising; Simon Dietz; Ivan Rudik; Gernot Wagner. 2024. \emph{Synthesis of evidence yields high social cost of carbon due to structural model variation and uncertainties}. Introduction; Section 1.1; Section 2.
\url{https://pmc.ncbi.nlm.nih.gov/articles/PMC11670083/}.
DOI: \url{https://doi.org/10.1073/pnas.2410733121}.
\item\hypertarget{definitions:source:30007714:2}{} Stefano Giglio; Theresa Kuchler; Johannes Stroebel; Olivier Wang. 2024. \emph{The Economics of Biodiversity Loss}. Working-paper record and abstract.
\url{https://www.nber.org/papers/w32678}.
DOI: \url{https://doi.org/10.3386/w32678}.
\end{enumerate}

\subsection*{Common conventions: Source mathematical conventions}

These conventions apply to each entry unless explicitly replaced.

\textbf{Models and identification.} A primitive model \(m\) specifies all probability laws, feasible choices, information, equilibrium and measurement rules used by its target. The admissible class \(\mathcal M\) is fixed independently of outcomes. If \(P_O(m)\) is its observable probability law and \(\tau(m)\) the target, its identified set at observable law \(P\) is
\[
 \mathcal I_\tau(P)=\{\tau(m):m\in\mathcal M,\ P_O(m)=P\}.
\]
Point identification means that this set is a singleton. An empty set signals incompatibility of data and assumptions. Coordinate bounds need not describe the joint set. Bounds are sharp only if every retained target is attainable or arbitrarily approachable by compatible models. Finite bounds require explicit restrictions on unobserved quantities; unrestricted confounding, hidden wealth, extrapolation or arbitrary equilibrium selection can yield uninformative sets.

\textbf{Causal interventions.} A potential outcome \(Y_i(p)\) is the outcome under a complete policy regime \(p\), including timing, financing, information and market spillovers. The target population and units are fixed before treatment. With interference, use the complete assignment vector or a justified exposure mapping; individual no-interference notation alone cannot identify an economy-wide intervention. A difference of mean potential outcomes does not require the cross-world joint distribution. Conditioning on post-treatment migration, survival, enrollment or employment changes the estimand and needs separate justification.

\textbf{Statistical statements.} An estimator, confidence set, forecast or test specifies a sampling law, model class and loss/coverage criterion. Uniform \((1-\alpha)\)-coverage means \(\inf_{m\in\mathcal M}\Pr_m\{\tau(m)\in C_n\}\geq1-\alpha\), or an explicitly asymptotic version. The whole parameter space trivially has coverage, so informativeness requires a stated width, risk or power objective. A forecasting improvement is relative to a named benchmark, information set and evaluation population.

\textbf{Equilibrium and optimization.} An equilibrium comprises feasible plans, optimality, beliefs where needed, budget constraints and all market/resource clearing. The primitives or a stated benchmark define each correspondence \(\mathcal E(p;m)\). Nonemptiness is assumed or proved, never inferred just from compact policy sets. A selected-equilibrium target uses a declared selection rule; a robust target quantifies over all permitted equilibria. Use suprema or infima unless attainment is established. An \(\varepsilon\)-optimal policy is feasible and within \(\varepsilon\) of the finite optimum value. Report an empty feasible set separately.

\textbf{Welfare and accounting.} Normative utility, interpersonal normalization, social weights, numeraire, discounting and the use of public receipts are primitives. Behavioral utility need not coincide with the welfare criterion. Household welfare includes ownership income and rents once; adding profits again to compensating variation generally double-counts them. A monetary equivalent is defined by an explicit transfer and price convention, with monotonicity and range conditions. Average effects, mechanism contributions, transfers and welfare losses are different targets. Infinite sums require convergence and dynamic feasible plans require terminal or no-Ponzi conditions.

\textbf{Source versions.} A working-paper date, revision date and journal publication year are distinguished. A DOI for a journal version is not automatically the DOI of an earlier manuscript. Locators refer to the stated version and printed pages where specified. A metadata-only check does not verify a claim about a full-text conclusion. Every entry's source subsection narrows the provenance claim accordingly.
% END ECONOMICS STATEMENT
\end{document}
